% !TeX root = ../navier-stokes-manuscript.tex
% ============================================================
% 3.2 The Compatible Intersection
% ============================================================
\subsection{The Compatible Intersection}\label{sub:TheCompatibleIntersection}

Every finite pair now has a terminal record.  Enlarging either cutoff retains
the lower record verbatim inside the larger one, because shared entries are the
same distributional derivatives of \(u\).  Order finite pairs coordinatewise
and assemble them with the bonding maps into
\[
  \mathscr X_\infty(I_*)
  :=
  \varprojlim_{N,M}\mathscr X_{N,M}(I_*).
\]
A family belongs to this inverse limit precisely when every finite coordinate
is an admissible record and
\[
  \rho_{N',M';I_*}^{N,M;I_*}(\mathbf R_{N,M})
  =\mathbf R_{N',M'}
\]
for every lower pair \((N',M')\).

\begin{proposition}[Compatible projective assembly]\label{prop:CompatibleProjectiveAssembly}
The datum-generated rectangles define
\[
  \cI[u_0]
  :=
  \bigl(\cR_{N,M}[u_0]\bigr)_{N,M}.
\]
This family lies in \(\mathscr X_\infty(I_*)\).  Its velocity coordinate is the
same distribution \(u\) throughout and obeys
\begin{equation}\label{eq:SectionThreeMixedAmbientSpace}
  u\in
  \bigcap_{r=0}^{\infty}
  \bigcap_{m=0}^{\infty}
  H^r(I_*;H^m(\R^3)).
\end{equation}
Coordinate projections commute with all bonding maps.
\end{proposition}

\begin{proof}
Given finite \(r,m\), choose \(N\geq r\) and \(M\geq m\).  That rectangle
contains
\[
  \partial_t^n u\in L^2(I_*;H^{M+1})
  \hookrightarrow L^2(I_*;H^m)
  \qquad(0\leq n\leq r).
\]
Because these are the distributional time derivatives of one velocity,
\(u\in H^r(I_*;H^m)\).  The pair was arbitrary, proving
\eqref{eq:SectionThreeMixedAmbientSpace}.  Compatibility makes each lower
record the literal projection of every larger one, so no representatives have
to be selected or reconciled.
\end{proof}

Project now to the zeroth temporal row while allowing every finite spatial
height:
\begin{equation}\label{eq:SectionThreeCompleteSpatialRow}
  \Szero(I_*)
  :=\bigcap_{m=0}^{\infty}L^2(I_*;H^m(\R^3)).
\end{equation}
Each \(\cR_{0,m}[u_0]\) supplies one coordinate, and compatibility identifies
them all with the same \(u\).  Hence
\[
  \pi_0\bigl(\cI[u_0]\bigr)
  =u\in\Szero(I_*).
\]

Every finite spatial energy used in the derivative ascent is a coordinate of
this row.  For \(s\geq0\), set
\begin{equation}\label{eq:SectionThreeSpatialEnergy}
  E_s(t):=\frac12\norm{\Lambda^su(t)}_2^2.
\end{equation}
Choose one integer \(m\geq s\).  Its coordinate gives
\begin{equation}\label{eq:SectionThreeAllSpatialEnergies}
  \int_0^{T_*}E_s(t)\,dt
  \leq
  C_{s,m}\norm{u}_{L^2(I_*;H^m)}^2
  <\infty.
\end{equation}
Each finite energy is therefore integrable on \(I_*\), with its own bound and
the same underlying velocity.
